\section{Conclusion}
\label{sec:conclusion}

Long-run population projection turns on a single quantity nobody can measure:
what fertility does after the demographic transition. Conventional practice
handles this by modelling the national rate as a mean-reverting time series,
which is a defensible statistical choice and an odd substantive one, because it
assumes away the fact that the parents of each generation are sampled in
proportion to how many children they have.

Adding that compositional force to a cohort-component projection shows it to be
real, measurable from evidence independent of the series it explains, and
consequential: \LadderMainstreamGain{} billion people at \ProjectionEndYear{},
or \LadderMainstreamGainPercent{}, from two parameters that statistical offices
already have the data to estimate. It also shows it to be more modest than either
its advocates or its critics would suggest. It shifts the level of the long-run
curve, it does not restore growth, and it is almost entirely a property of
ordinary variation rather than of the visible minorities that dominate discussion
of the subject.

Against that measured force sits an unmeasurable one, and the paper's main
contribution is a way of reporting the comparison that does not require guessing.
Both act multiplicatively on fertility, so ``what will the population be'' can be
replaced by ``how fast would conditions have to keep deteriorating for the
compositional force to lose'' --- \BoundaryPaired{} per decade to cancel its
effect on the fertility rate, \BoundaryPopulation{} to cancel its effect on the
population at the horizon, and \BoundaryFloor{} to \BoundaryCeiling{} for the
first of those across the empirical uncertainty in the two measured parameters.

We do not know which side of that threshold the world is on. We would rather
report the threshold than pretend otherwise, and rather be wrong for a reason
that can be pointed at on \Cref{fig:boundary} than be approximately right for
none.
